% Independent supporting proof fragment. No standalone preamble.
% This is a reduction to established tools, not a literature-priority claim.
\section{Finite-transform reduction of a residual transfer certificate}
\label{supp:abstract-transfer-reduction}

I state the mechanism with coordinate-specific positive measures so that
its dependence on Gaussianity is explicit. Let $I_j$ be a compact
interval, let $\nu_j$ be a finite positive measure on $I_j$, and partition
$I_j$ into cells $I_{jc}$ with actual nodes $c_{jc}\in I_{jc}$. Set
$\tau_j=\sum_c\nu_j(I_{jc})\delta_{c_{jc}}$. Suppose positive finite
measures $\rho_{j1},\rho_{j2}$ satisfy
\begin{equation}
 |(\nu_j-\tau_j)(f)|\leq
 \rho_{j1}(|f'|)+\rho_{j2}(|f''|),\qquad f\in C^2(I_j).
 \label{eq:abstract-coordinate-bound}
\end{equation}
For the cell construction in the project, this follows from the
established centred Taylor identity. It also holds for an arbitrary
finite positive cell measure: the second kernel has density
$\int_{I_{jc}\cap(-\infty,z]}(z-y)\nu_j(dy)$ below its node and
$\int_{I_{jc}\cap[z,\infty)}(y-z)\nu_j(dy)$ above it. The first
measure has atoms dominating
$|\int_{I_{jc}}(y-c_{jc})\nu_j(dy)|$. An atom of $\nu_j$ at a shared
cell boundary is assigned to exactly one cell. This convention avoids
double counting. Gaussian equal-CDF cells are sufficient, rather than
necessary, assumptions for \eqref{eq:abstract-coordinate-bound}.

For $X_i(y)=e^{d_i+s_i\cdot y}$, let $\mathcal A_D$ be the finite
exponential-polynomial space spanned by $X^\alpha$ with
$\alpha\in\mathbb N_0^n$ and
$|\alpha|\leq D$. Signed coefficients are permitted. For
$1\leq j\leq p$ and $k\in\{1,2\}$, define
\[
 \Lambda_{jk}=\nu_1\otimes\cdots\otimes\nu_{j-1}
       \otimes\rho_{jk}\otimes\tau_{j+1}\otimes\cdots\otimes\tau_p.
\]
Suppose $R$ is measurable on the full domain, belongs to
$C^2(\prod_jI_j)$ on the retained box, and has pointwise derivative
envelopes
\[
 |\partial_j^kR|\leq E_{jk}
 =\sum_\ell a_{jk\ell}H_{jk\ell}
   +\sum_q b_{jkq}|A_{jkq}B_{jkq}|,
 \qquad a_{jk\ell},b_{jkq}\geq0.
\]
Each $H$ is nonnegative on the retained box. Each $H,A^2,B^2$ has a
known finite representation in $\mathcal A_D$. Whenever several pairs
$(A_h,B_h)$ satisfy the same pointwise identity
$|A_hB_h|=|AB|$, any of their Cauchy bounds may be used and their
minimum is valid. Finally suppose an explicit finite bound $T_{\rm out}$
controls the integral of $|R|$ outside the retained box under a full
product law $\mu_1\otimes\cdots\otimes\mu_p$, whose restriction is
$\nu_1\otimes\cdots\otimes\nu_p$.

\begin{proposition}[A finite-transform transfer reduction]
\label{prop:abstract-transfer-reduction}
Under these assumptions,
\[
 \left|\int R\,d(\mu_1\otimes\cdots\otimes\mu_p)
       -\int R\,d(\tau_1\otimes\cdots\otimes\tau_p)\right|
 \leq T_{\rm out}+\sum_{j=1}^p\sum_{k=1}^2 \mathcal B_{jk},
\]
where
\[
 \mathcal B_{jk}=
 \sum_\ell a_{jk\ell}\Lambda_{jk}(H_{jk\ell})
 +\sum_q b_{jkq}
     \sqrt{\Lambda_{jk}(A_{jkq}^2)\Lambda_{jk}(B_{jkq}^2)}.
\]
The stated alternative-pair minima may replace the corresponding
square-root terms. Every moment is a finite sum of products of scalar
Laplace transforms. In particular, for
$h=\sum_\alpha b_\alpha X^\alpha$,
\begin{equation}
 \Lambda_{jk}(h)=\sum_\alpha b_\alpha
   e^{\sum_i\alpha_id_i}
   \mathcal M_{\rho_{jk}}\!\left(\sum_i\alpha_is_{ij}\right)
   \prod_{\ell<j}\mathcal M_{\nu_\ell}\!\left(\sum_i\alpha_is_{i\ell}\right)
   \prod_{\ell>j}\mathcal M_{\tau_\ell}\!\left(\sum_i\alpha_is_{i\ell}\right).
 \label{eq:abstract-factorised-moments}
\end{equation}
Thus certified scalar transforms at the finite set of required rates
suffice to construct a price-independent transfer allowance. No value of
$R$ at a Cartesian grid point is required to construct this bound.
\end{proposition}
\begin{proof}
The product difference telescopes exactly as
\[
 \bigotimes_j\nu_j-\bigotimes_j\tau_j
 =\sum_j\left(\bigotimes_{\ell<j}\nu_\ell\right)
          \otimes(\nu_j-\tau_j)
          \otimes\left(\bigotimes_{\ell>j}\tau_\ell\right).
\]
For fixed other coordinates, apply
\eqref{eq:abstract-coordinate-bound} to $R$ as a function of its
$j$th coordinate. Compact support and $C^2$ regularity imply bounded
derivatives and justify all applications of Fubini. Integrating the
coordinate inequality against the other positive measures and summing
gives $\sum_{jk}\Lambda_{jk}(|\partial_j^kR|)$. The pointwise envelope
and Cauchy's inequality in $L^2(\Lambda_{jk})$ give the displayed
$\mathcal B_{jk}$. Positivity, rather than probability normalisation,
is the required property: no total-mass divisor occurs. The outside
integral contributes at most $T_{\rm out}$ by the triangle inequality.

For a monomial, $X^\alpha=e^{\sum_i\alpha_id_i}
\prod_\ell e^{(\sum_i\alpha_is_{i\ell})y_\ell}$.
Its integral under the product measure therefore factorises. Linearity
of a finite sum gives \eqref{eq:abstract-factorised-moments}, including
its signed coefficients. At zero rate each transform is the actual
measure mass. This proves both the identity and its mass convention.
Certified interval arithmetic preserves the enclosure when signed
products are summed. A known nonnegative moment permits intersecting
its interval with $[0,\infty)$, but a negative upper endpoint signals
an invalid enclosure rather than permission to substitute zero.
\end{proof}

\subsection{Why the present residual needs only degree four}

For positive basket weights of sum one, weighted-centred drift and
loadings give
$T=\sum_iw_iX_i\geq\exp(\sum_iw_i(d_i+s_i\cdot y))=1$.
Hence $U=T-1\geq0$. With $u=T_j$, $v=T_{jj}\geq0$ and affine
$m=m_G+r\cdot y$, the exact conditional second residual is
\[
 R_2=U^2\int_0^1(1-z)t^{-2}C_2(m+\log t)dz,
 \qquad t=1+zU\geq1.
\]
If $C_3=C_2'-2C_2$ and $C_4=C_3'-3C_3$, differentiation of
$G=t^{-2}C_2(m+\log t)$ gives
\[
 G_j=r_jt^{-2}C_2'+zu t^{-3}C_3,
\]
\[
 G_{jj}=r_j^2t^{-2}C_2''+2r_jzu t^{-3}C_3'
                         +zv t^{-3}C_3+z^2u^2t^{-4}C_4.
\]
The derivatives of $U^2G$ are $2UuG+U^2G_j$ and
$2(u^2+Uv)G+4UuG_j+U^2G_{jj}$. These expressions exhibit exactly
which bounded boundary derivatives are required. Uniform bounds on
$C_2,C_3,C_4,C_2',C_2'',C_3'$ and integration of $1-z$,
$z(1-z)$ and $z^2(1-z)$ produce the project's two displayed
envelopes. This residual differentiation is the extra argument that a
generic Peano identity does not supply.

The direct nonnegative factors are
$U^2,u^2,Uv,Uu^2,U^2v,U^2u^2$. Their degrees in $X$ are at most
four. The only absolute products needing a closure are $U|u|$ and
$U^2|u|$. Cauchy bounds the first using $U^2,u^2$, and the second
using either $U^2,U^2u^2$ or $U^4,u^2$. Consequently all required
moments have degree at most four. Neither Cauchy nor exponential
factorisation needs a new principle once these pointwise derivative
envelopes have been obtained.

For a fixed finite number of payoff kinks and a strictly positive
independent Gaussian smoothing scale, the falling log derivatives are
finite sums of polynomial multiples of Gaussian boundary densities.
Each is bounded, so this argument includes the stated vanilla and
capped calls. The scale-dependent constants are not uniform as the
smoothing scale tends to zero. This boundedness statement concerns the
falling boundary derivatives; the vanilla kernel and its delta remain
unbounded. A general multivariate payoff needs its own smoothness and
derivative envelopes and does not inherit this degree-four conclusion.

\subsection{Construction cost and what the reduction does not prove}

For fixed degree $D$, the number of distinct monomials satisfies
$M\leq\binom{n+D}{D}$. For a fixed number of geometry polynomials,
forming the other-coordinate products naively costs $O(p^2M)$
arithmetic products. Grouping by the transported-coordinate rate and
reusing those products gives $O(BpM)$ subsequent dot products for $B$
blocks, in addition to rate formation, coefficient construction, scalar
transform evaluation and certified rounding. A cached prefix or suffix
scheme could reduce the other-coordinate product count, but that would
be an implementation improvement requiring its own verified producer;
it is not credited to the present implementation.

The rounded scalar transforms still involve $N$ actual nodes. Direct
evaluation costs $O(N)$ scalar terms per query. A factorial-series
prefix implementation pays $O(NJ)$ preparation at series order $J$,
then $O(J)$ arithmetic operations per block/rate query, plus continuous
transform evaluations. Its remainder for a signed finite measure
$\eta$ supported in $[-L,L]$ is at most
$\|\eta\|_{\rm TV}e^{|v|L}(|v|L)^{J+1}/(J+1)!$.
These are arithmetic counts, not a complete uniform bit-complexity
theorem. At $b$ bits, each count must be weighted by the costs of the
relevant arithmetic and transcendental operations; node inversion,
exact rational rates, output magnitude and cancellation can require
additional precision. Fixed $D$ is essential: variable $D$ makes the
monomial count combinatorial. None of these statements is a
dimension-free error guarantee or a quantum pricing advantage.

Gaussianity supplies explicit scalar transforms and the exact tilted
outside-box moment used in this project. It is not required by
Proposition~\ref{prop:abstract-transfer-reduction}. Independent
non-Gaussian coordinates with certified finite transforms and an
explicit outside allowance also satisfy the abstract proposition.
Correlated coordinates require either a justified independent latent
representation or a joint-transform construction. Their marginal
transforms alone do not determine the required moments.

The displayed proposition is an independently derived packaging of
established Taylor transport, product telescoping, Cauchy and
exponential moment factorisation. Its mathematical correctness does not
establish literature originality. The project-specific conclusion is
the applicability of this package to the stated residual, including its
verified derivative constants, actual rounded law and observed monetary
allowances. A material claim beyond this applicability still needs a
source-level comparison and competitive evidence.
